\documentclass[11pt,a4paper,sans]{moderncv}
\moderncvstyle{classic}
\moderncvcolor{blue}
\usepackage[utf8]{inputenc}
\usepackage[scale=0.93]{geometry} % scale=0.93 (contre 0.8 pour les CV) permet de tenir sur une page

\name{Alexandro}{Disla}
\title{Lettre de Motivation}
\address{Rue Saint-Louis Jeanty \#14, Route de Frère}{Port-au-Prince, Haïti}{}
\phone[mobile]{(+509) 4148-3700}
\email{alexandrodisla@hotmail.com}
\extrainfo{GitHub: \url{https://github.com/AD0791}}

\begin{document}

\recipient{Équipe de Recrutement de Samaritan's Purse Haïti}{Bureau Pays d'Haïti\\Poste : Senior MEAL Officer}
\date{\today}
\opening{Madame, Monsieur,}
\closing{Je vous prie d'agréer, Madame, Monsieur, l'expression de mes salutations distinguées,}
\enclosure[Ci-joint]{Curriculum Vitae}

\makelettertitle

C'est avec un vif intérêt que je vous adresse ma candidature au poste de Senior MEAL Officer au sein du bureau pays de Samaritan's Purse en Haïti. J'apporte plus de cinq ans d'expérience MEAL sur trois secteurs dans ce pays --- la santé à la Caris Foundation International, l'eau et l'assainissement chez HANWASH, l'éducation chez Anseye Pou Ayiti --- doublés de la capacité d'ingénierie qui permet de construire et de sécuriser les systèmes de données sur lesquels ces programmes reposent. Je réside à Port-au-Prince, je travaille quotidiennement en français, en anglais et en créole haïtien, et je suis disponible pour les déplacements de terrain et les visites de vérification.

Le cœur de ce poste, tel que je le lis, consiste à garantir que les données rapportées par un bureau pays résistent à la vérification lorsqu'on les remonte jusqu'à leur source. C'est précisément le travail que j'ai fait. Chez HANWASH, j'ai raffiné les plans MEAL et les tableaux de suivi des indicateurs alignés sur les normes JMP pour des programmes d'eau et d'assainissement. À la Caris Foundation, j'ai établi et supervisé les protocoles d'assurance qualité des données (DQA) sur des données intégrées MySQL, en remontant les chiffres rapportés jusqu'aux registres sources et en suivant les actions correctives jusqu'à leur clôture, et non jusqu'à leur simple consignation. Vérifier à la source et refermer la boucle : c'est la discipline que j'apporterais au cadre MEAL du bureau pays et aux visites de vérification terrain.

J'aborde la protection des données comme un problème d'ingénierie plutôt que comme un document de politique. Sur CommCare, mWater, ODK et KoboToolbox, j'ai conçu des instruments de collecte avec logiques de saut et contraintes de validation afin que les erreurs soient arrêtées à la saisie, et j'ai administré les bases qui les sous-tendent : contrôle d'accès par rôle, stockage cloud sécurisé sur Azure et AWS, dédoublonnage systématique des listes de bénéficiaires. La confidentialité et la durée de conservation se décident dans le schéma et dans les permissions ; je suis à l'aise pour les décider et les documenter.

En matière de reporting et d'appui à la décision, j'ai construit et maintenu des tableaux de bord programme sous Power BI, Quarto et mWater pour un suivi hebdomadaire des indicateurs, et automatisé le reporting d'activité en Python et SQL à la Caris Foundation, réduisant le temps de traitement manuel de 40 \%. Économiste appliqué de formation, je sais mener une analyse de l'échantillonnage jusqu'à la tendance, et --- ce qui compte tout autant --- énoncer clairement à un chef de programme ou à un Deputy Country Director ce que cette tendance signifie, ce qui y demeure incertain et quelle décision elle appelle.

J'ai supervisé des agents de collecte sur des programmes multi-sites, assuré l'encadrement technique de personnel de suivi-évaluation, formé des équipes de terrain aux outils numériques et aux tableaux de bord, et coordonné avec des partenaires gouvernementaux dont la DINEPA et le MSPP. Sur la redevabilité envers les populations affectées, je souhaiterais construire le mécanisme de retour et de plainte comme tout autre système de données : un canal que les personnes peuvent réellement atteindre, un enregistrement confidentiel, des voies de référencement définies à l'avance et un suivi jusqu'à la clôture --- afin que ce que les communautés nous disent modifie le programme, au lieu de finir dans un classeur.

Ce poste m'attire parce que Samaritan's Purse sert les communautés haïtiennes sans interruption depuis 2010, et parce qu'il traite les données probantes et la redevabilité comme des obligations envers les personnes servies plutôt que comme une charge de reporting. Je respecte la mission et les valeurs chrétiennes de l'organisation ainsi que son engagement auprès des populations affectées par les crises, et je serais heureux de contribuer à un bureau pays où cet engagement façonne la manière dont le travail est mené. Je reste à votre entière disposition pour un entretien.

\makeletterclosing

\end{document}
